% Part: first-order-logic
% Chapter: sequent-calculus
% Section: proof-theoretic-notions

\documentclass[../../../include/open-logic-section]{subfiles}

\begin{document}

\begin{editorial}
  این بخش تعریف‌های رابطهٔ اثبات‌پذیری و سازگاری را برای حساب دنباله‌ای
  گرد می‌آورد.
\end{editorial}

\iftag{FOL}
      {\olfileid{fol}{seq}{ptn}}
      {\olfileid{pl}{seq}{ptn}}

\olsection{مفاهیم برهان‌نظری}

\begin{explain}
همان‌گونه که چند مفهوم معناشناختی مهم (اعتبار، استلزام و ارضاپذیری)
را تعریف کردیم، اکنون \emph{مفاهیم برهان‌نظریِ} متناظر را تعریف
می‌کنیم. این مفاهیم نه با توسل به ارضای !!{sentence}s در
\iftag{FOL}{!!{structure}s}{!!{valuation}s}، بلکه با توسل به
!!{derivability} یا !!{nonderivability} برخی دنباله‌ها تعریف می‌شوند.
کشف انطباق این مفاهیم اهمیتی بسیار داشت. این انطباق محتوای قضیه‌های
\emph{صحت} و \emph{تمامیت} است.
\end{explain}

\begin{defn}[قضایا]
!!^a{sentence}~$!A$ \emph{قضیه} است هرگاه در~$\Log{LK}$،
!!a{derivation}ی برای دنبالهٔ $\quad \Sequent !A$ وجود داشته باشد.
$\Proves !A$ می‌نویسیم اگر $!A$ قضیه باشد، و اگر قضیه نباشد
$\Proves/ !A$ می‌نویسیم.
\end{defn}

\begin{defn}[!!^{derivability}]
برای !!^a{sentence}~$!A$، عبارت \emph{!!{derivable} از مجموعه‌ای از}
!!{sentence}s چون~$\Gamma$ را به کار می‌بریم، و
$\Gamma \Proves !A$ می‌نویسیم، اگر و تنها اگر
زیرمجموعهٔ متناهی~$\Gamma_0 \subseteq \Gamma$ و توالی‌ای چون
$\Gamma_0'$ از !!{sentence}s موجود در~$\Gamma_0$ وجود داشته باشد
به‌طوری‌که $\Log{LK}$، $\Gamma_0' \Sequent !A$ را !!{derive}s کند.
اگر $!A$ از $\Gamma$ !!{derivable} نباشد، $\Gamma \Proves/ !A$
می‌نویسیم.
\end{defn}

به سبب قواعد انقباض، تضعیف و تبادل، ترتیب و شمار !!{sentence}s در
~$\Gamma_0'$ اهمیتی ندارد: اگر دنبالهٔ
$\Gamma_0' \Sequent !A$ !!{derivable} باشد، آنگاه
$\Gamma_0'' \Sequent !A$ نیز برای هر $\Gamma_0''$ که همان
!!{sentence}s موجود در~$\Gamma_0'$ را دربر دارد چنین است.
برای نمونه، اگر $\Gamma_0 = \{!B, !C\}$ باشد، آنگاه هم
$\Gamma_0' = \tuple{!B, !B, !C}$ و هم
$\Gamma_0'' = \tuple{!C, !C, !B}$ توالی‌هایی‌اند که دقیقاً
!!{sentence}s موجود در~$\Gamma_0$ را دربر دارند. اگر دنباله‌ای که
یکی از آن‌ها را دربر دارد !!{derivable} باشد، دنبالهٔ دیگر نیز چنین
است؛ برای نمونه:
\begin{prooftree}
  \AxiomC{}
  \Deduce$!B, !B, !C \fCenter !A$
  \RightLabel{\LeftR{\Contraction}}
  \UnaryInf$!B, !C \fCenter !A$
  \RightLabel{\LeftR{\Exchange}}
  \UnaryInf$!C, !B \fCenter !A$
  \RightLabel{\LeftR{\Weakening}}
  \UnaryInf$!C, !C, !B \fCenter !A$
\end{prooftree}
از این پس می‌گوییم اگر $\Gamma_0$ مجموعه‌ای متناهی از !!{sentence}s
باشد، آنگاه $\Gamma_0 \Sequent !A$ هر دنباله‌ای است که مقدمش توالی‌ای
از !!{sentence}s موجود در~$\Gamma_0$ باشد؛ و در صورت لزوم انقباض‌ها،
تبادل‌ها و تضعیف‌ها را به‌طور ضمنی منظور می‌کنیم.

\begin{defn}[سازگاری]
مجموعه‌ای از جمله‌ها چون~$\Gamma$ \emph{ناسازگار} است اگر و تنها اگر
زیرمجموعهٔ متناهی~$\Gamma_0 \subseteq \Gamma$ وجود داشته باشد
به‌طوری‌که $\Log{LK}$، $\Gamma_0 \Sequent \quad$ را !!{derive}s کند.
اگر $\Gamma$ ناسازگار نباشد، یعنی اگر برای هر
$\Gamma_0 \subseteq \Gamma$ متناهی، $\Log{LK}$،
$\Gamma_0 \Sequent \quad$ را !!{derive} نکند، می‌گوییم
\emph{سازگار} است.
\end{defn}

\begin{prop}[بازتابی]
\ollabel{prop:reflexivity}
اگر $!A \in \Gamma$، آنگاه $\Gamma \Proves !A$.
\end{prop}

\begin{proof}
دنبالهٔ آغازینِ $!A \Sequent !A$ !!{derivable} است و
$\{!A\} \subseteq \Gamma$.
\end{proof}

\begin{prop}[یکنوایی]
\ollabel{prop:monotonicity}
اگر $\Gamma \subseteq \Delta$ و $\Gamma \Proves !A$، آنگاه
$\Delta \Proves !A$.
\end{prop}

\begin{proof}
فرض کنید $\Gamma \Proves !A$؛ یعنی $\Gamma_0 \subseteq \Gamma$
متناهی‌ای وجود دارد که $\Gamma_0 \Sequent !A$ !!{derivable} است.
چون $\Gamma \subseteq \Delta$، پس $\Gamma_0$ زیرمجموعه‌ای متناهی
از~$\Delta$ نیز هست. بنابراین !!{derivation} مربوط به
$\Gamma_0 \Sequent !A$ همچنین $\Delta \Proves !A$ را نشان می‌دهد.
\end{proof}

\begin{prop}[تعدی]
\ollabel{prop:transitivity}
اگر $\Gamma \Proves !A$ و $\{!A\} \cup \Delta \Proves !B$، آنگاه
$\Gamma \cup \Delta \Proves !B$.
\end{prop}

\begin{proof}
اگر $\Gamma \Proves !A$، آنگاه مجموعهٔ متناهی
$\Gamma_0 \subseteq \Gamma$ و !!a{derivation}~$\pi_0$ برای
$\Gamma_0 \Sequent !A$ وجود دارند. اگر
$\{!A\} \cup \Delta \Proves !B$، آنگاه برای زیرمجموعهٔ متناهی‌ای چون
$\Delta_0 \subseteq \Delta$، !!a{derivation}~$\pi_1$ برای
$!A, \Delta_0 \Sequent !B$ وجود دارد. !!{derivation} زیر را در نظر
بگیرید:
\begin{prooftree}
  \AxiomC{}
  \RightLabel{$\pi_0$}
  \Deduce$\Gamma_0 \fCenter !A$
  \AxiomC{}
  \RightLabel{$\pi_1$}
  \Deduce$!A, \Delta_0 \fCenter !B$
  \RightLabel{\Cut}
  \BinaryInf$\Gamma_0, \Delta_0 \fCenter !B$
\end{prooftree}
چون $\Gamma_0 \cup \Delta_0 \subseteq \Gamma \cup \Delta$، این
نشان می‌دهد $\Gamma \cup \Delta \Proves !B$.
\end{proof}

توجه کنید که این به‌ویژه بدان معناست که اگر $\Gamma \Proves !A$ و
$!A \Proves !B$، آنگاه $\Gamma \Proves !B$. همچنین نتیجه می‌شود که
اگر $!A_1, \dots, !A_n \Proves !B$ و $\Gamma \Proves !A_i$ برای
هر~$i$، آنگاه $\Gamma \Proves !B$.

\begin{prop}
\ollabel{prop:incons}
$\Gamma$ ناسازگار است اگر و تنها اگر $\Gamma \Proves {!A}$ برای هر
جملهٔ~$!A$.
\end{prop}

\begin{proof}
تمرین.
\end{proof}

\begin{prob}
\olref[fol][seq][ptn]{prop:incons} را ثابت کنید.
\end{prob}

\begin{prop}[فشردگی]
  \ollabel{prop:proves-compact}
  \begin{enumerate}
  \item اگر $\Gamma \Proves !A$، آنگاه زیرمجموعهٔ متناهی
    $\Gamma_0 \subseteq \Gamma$ وجود دارد که $\Gamma_0 \Proves !A$.
  \item اگر هر زیرمجموعهٔ متناهیِ~$\Gamma$ سازگار باشد، آنگاه
    $\Gamma$ سازگار است.
  \end{enumerate}
\end{prop}

\begin{proof}
  \begin{enumerate}
    \item اگر $\Gamma \Proves !A$، آنگاه زیرمجموعهٔ متناهی
      $\Gamma_0 \subseteq \Gamma$ وجود دارد که دنبالهٔ
      $\Gamma_0 \Sequent !A$ دارای !!a{derivation} است. در نتیجه،
      $\Gamma_0 \Proves !A$.
    \item اگر $\Gamma$ ناسازگار باشد، زیرمجموعهٔ متناهی
      ~$\Gamma_0 \subseteq \Gamma$ وجود دارد که $\Log{LK}$،
      $\Gamma_0 \Sequent \quad$ را !!{derive}s می‌کند. اما در این صورت،
      $\Gamma_0$ زیرمجموعه‌ای متناهی و ناسازگار از~$\Gamma$ است.
  \end{enumerate}
\end{proof}

\end{document}
